\section{The canonical predictor}\label{sec:predictor}

In the exact count, each placement is a list of occupied squares. The predictor uses a different variable:
a weight $p_{ij}$ at each square. The row, column and diagonal sums express the constraints retained by the
model.
Among all admissible assignments of weights, we choose just one: the maximizer of the functional below.

\paragraph{The canonical functional.}
Let $\mathcal D$ denote the union of the two diagonal families, with $|\mathcal D|=4n-2$. Consider the
polytope and the functional
\begin{equation}\label{eq:Kn}
 K_n=\Bigl\{p\in[0,1/4]^{n^2}:\ \sum_jp_{ij}=\sum_ip_{ij}=1,\ m_d(p)\le1\ (d\in\mathcal D)\Bigr\},
\end{equation}
\begin{equation}\label{eq:Hn}
 H_n(p)=\sum_{ij}\bigl\{-p_{ij}\log p_{ij}+3(1-p_{ij})\log(1-p_{ij})\bigr\}-\sum_{d\in\mathcal D}h_d(p)\log
h_d(p),
\end{equation}
with $0\log0=0$. This functional is the Bethe entropy~\cite{YFW} of the system of rows, columns and
diagonals, expressed in the square variables and restricted to fractional placements of total mass $n$.
The cap $p\le1/4$ makes the square term concave.

\begin{lema}[The canonical maximizer]\label{lem:maximizador}
For $n\ge4$, $K_n$ is a nonempty compact convex set, and $H_n$ is continuous on $K_n$ and strictly concave
along every segment in $K_n$. Consequently, $H_n$ has a unique maximizer on $K_n$, which we
denote by $\pi^{(n)}$ and call the canonical maximizer.
\end{lema}
\begin{proof}
The uniform placement $p\equiv1/n$ lies in $K_n$ if $n\ge4$, because each diagonal has at most $n$
squares. We write $H_n(p)=\sum_cf(p_c)+\sum_{d\in\mathcal D}G(m_d(p))$, with
$f(t)=-t\log t+3(1-t)\log(1-t)$ on $[0,1/4]$ and $G(z)=-(1-z)\log(1-z)$ on $[0,1]$. Both functions are
continuous on
their closed intervals, and $f''(t)=-1/t+3/(1-t)<0$ on $(0,1/4)$.

Let $p^0\ne p^1$ in $K_n$, $p^t=(1-t)p^0+tp^1$ and $I=\{c:p^0_c\ne p^1_c\}\ne\emptyset$. For $c\in I$ and
$t\in(0,1)$, $p^t_c$ lies strictly between $p^0_c$ and $p^1_c$, hence in $(0,1/4)$; thus
$t\mapsto f(p^t_c)$ is of class $C^2$ on $(0,1)$ with second derivative $f''(p^t_c)(p^1_c-p^0_c)^2<0$.
Squares outside $I$ give constants. Each $t\mapsto G(m_d(p^t))$ is concave because $G$ is concave
on $[0,1]$ and $m_d$ is affine. The sum $\varphi(t)=H_n(p^t)$ is continuous on $[0,1]$ and strictly concave
on $(0,1)$, hence concave on $[0,1]$. If the strict inequality failed between two points
$t_0<t_1$ of $[0,1]$, $\varphi$ would touch its chord at an interior point of $[t_0,t_1]$ and would be
affine on
that interval, contrary to strict concavity on $(t_0,t_1)$. Existence of the maximizer follows
from compactness, and uniqueness from strict concavity.
\end{proof}

\begin{definicion}\label{def:Hinf}
We say that $(H_\infty)$ holds if the canonical maximizer is interior:
$0<\pi^{(n)}_c<1/4$ for every square $c$ and $h_d(\pi^{(n)})>0$ for every diagonal $d$.
\end{definicion}

Section~\ref{sec:P02} proves that $(H_\infty)$ holds for every $n\ge n_0$, with
$c/n\le\pi^{(n)}_{ij}\le M/n$ and $h_d\ge\eta$ for positive constants $c$, $M$ and $\eta$
(see~(P2.7)). The threshold $n_0$ given by that proof is very large (end of Section~\ref{sec:P02}).
For $n=4$ the hypothesis fails, and for $n=5$ the numerical computation indicates that it fails as well
(Section~\ref{sec:discusion}); for $6\le n<n_0$ it is not proved, and for the sizes of
Section~\ref{sec:numerico} it is checked numerically.
Below we write $\pi=\pi^{(n)}$ when there is no ambiguity.

\paragraph{The pseudodeterminant and the local term.}
Let $n\ge4$ be such that its maximizer satisfies $(H_\infty)$. Set
\begin{equation}\label{eq:gS}
 g_c=\frac{\pi_c(1-\pi_c)}{1-4\pi_c},\qquad
 S_n=\operatorname{diag}(0_R,0_C,h_D,h_A)+A^T\operatorname{diag}(g)A,
\end{equation}
a symmetric matrix indexed by the $6n-2$ lines, and set $v_0=(\mathbf1_R,-\mathbf1_C,0_D,0_A)$. The
quadratic form of $S_n$ is a sum of squares whose weights, $h_d$ and $g_c$, are positive by
$(H_\infty)$. If it vanishes at a vector
$x=(x_R,x_C,x_D,x_A)$, then $x_D=0$, $x_A=0$ and $x_{R,i}+x_{C,j}=0$ for each square $(i,j)$. Its
kernel is therefore the line spanned by $v_0$, and $S_n$ is positive definite on $v_0^\perp$. For
$n\ge n_0$, Section~\ref{sec:P03} also gives a uniform gap. We call the direction $v_0$, the discrete
analogue of the neutral mode $(1,-1,0,0)$
in~\cite[Section~7.2]{Perfil}, and its analogue on residual boards the \emph{gauge}:
since $Av_0=0$, adding a multiple of $v_0$ to a vector indexed by the lines does not change its
sums at the squares. Fixing the gauge means choosing a representative, for instance in $v_0^\perp$. The
pseudodeterminant
$\operatorname{pdet}S_n$ is the product of its nonzero eigenvalues. The local term is
\[
 L_n=-\frac12\sum_{ij}\bigl[\log(1-4\pi_{ij})-4\log(1-\pi_{ij})\bigr],
\]
as in~(LC1).

\begin{definicion}[Canonical predictor]\label{def:Qpred}
For $n\ge4$ satisfying $(H_\infty)$,
\begin{equation}\label{eq:Qpred}
 \log Q_{\rm pred}(n)=H_n(\pi^{(n)})+L_n-\tfrac12\log\operatorname{pdet}S_n+\tfrac12\log(4\pi n).
\end{equation}
\end{definicion}

\begin{observacion}[Origin of the definition]
Of the four terms in the definition, the first is the Bethe entropy at its maximum. The other
three correspond, heuristically, to two further ingredients: the one-loop correction to the
Bethe approximation, expressed as a determinant of the Hessian matrix by the
Ihara--Bass formula in the form of Watanabe and Fukumizu~\cite{Bass,WF}, and the Gaussian normalization of
the
zero mode that imposes equality between the number of empty rows and the number of empty columns. For these
approximations for permanents and systems on graphs, see~\cite{CC,ChertkovYedidia,Vontobel}, and for
loop corrections to the Bethe approximation,~\cite{MontanariRizzo}. In the rook model,
a formula with the same three ingredients approximates the logarithm of the count with error
$O(1/k)$ up to the full board (Appendix~\ref{sec:torres}). In this paper \eqref{eq:Qpred} is
only a definition: no proof uses this interpretation or the appendix.
\end{observacion}

For the residual boards that arise as queens are placed one at a time, the predictor $L(B)$ is defined in
Section~\ref{sec:B3T01}, with the same functional, the canonical determinant of the residual board and the
normalization $\tfrac12\log(2\pi)$ per component. On the full board, $L(B)$ coincides with
$\log Q_{\rm pred}(n)$ (Lemma~\ref{lem:puente}).
